\chapter{Software Requirements}

\section{User Interactions / Use Cases}
The system supports three primary user interactions:
\begin{enumerate}
    \item \textbf{Capture}: The user opens the Android app, starts recording, walks through the room while receiving live quality feedback, stops recording, and taps ``Upload \& Process.''
    \item \textbf{Monitor}: The user watches the processing progress screen showing current step, percentage, elapsed time, and log output.
    \item \textbf{View/Measure}: The user opens the finished 3D model in the web viewer and takes point-to-point measurements (FYP-2).
\end{enumerate}

\section{Use Case Diagram}
The system involves three actors (Figure~\ref{fig:usecase}): the \textit{End User} (captures and views), the \textit{Processing Server} (manages jobs), and the \textit{Cloud GPU} (runs reconstruction). The use cases are Record Scan, Upload Dataset, Monitor Progress, View Model, Take Measurement, and Configure Server. Each maps to code: Record Scan to \texttt{HelloArActivity} and \texttt{CaptureLogger}; Upload Dataset and Monitor Progress to \texttt{ProcessingActivity} and the \texttt{/upload} and \texttt{/status} endpoints of \texttt{server/main.py}; Configure Server to the app's settings menu (server URL) and \texttt{/config}; View Model and Take Measurement to the desktop viewer today and the web workspace in FYP-2.

% Stick-figure actor for the use case diagram: \bsactor{x}{y}{label}
\newcommand{\bsactor}[3]{%
  \begin{scope}[shift={(#1,#2)}]
    \draw (0,0.95) circle (0.2);
    \draw (0,0.75) -- (0,0.05);
    \draw (-0.35,0.5) -- (0.35,0.5);
    \draw (0,0.05) -- (-0.3,-0.45);
    \draw (0,0.05) -- (0.3,-0.45);
    \node[align=center] at (0,-0.85) {#3};
  \end{scope}}

\begin{figure}[htbp]
\centering
\begin{tikzpicture}[font=\footnotesize,
  uc/.style={draw, ellipse, minimum width=3.6cm, minimum height=0.85cm, align=center}]
  \draw[rounded corners=4pt] (-2.6,-3.6) rectangle (2.6,4.0);
  \node[font=\small\bfseries] at (0,3.6) {BuildScan 3D};
  \node[uc] (rec)  at (0,2.7)  {Record Scan};
  \node[uc] (upl)  at (0,1.6)  {Upload Dataset};
  \node[uc] (mon)  at (0,0.5)  {Monitor Progress};
  \node[uc] (view) at (0,-0.6) {View Model};
  \node[uc] (meas) at (0,-1.7) {Take Measurement};
  \node[uc] (conf) at (0,-2.8) {Configure Server};
  \bsactor{-5.2}{-0.2}{End User}
  \bsactor{5.2}{1.4}{Processing\\Server}
  \bsactor{5.2}{-2.4}{Cloud GPU}
  \coordinate (user) at (-4.8,0.3);
  \coordinate (srv)  at (4.8,1.9);
  \coordinate (gpu)  at (4.8,-1.9);
  \foreach \n in {rec,upl,mon,view,meas,conf} {\draw (user) -- (\n.west);}
  \foreach \n in {upl,mon,conf} {\draw (srv) -- (\n.east);}
  \draw (gpu) -- (upl.east);
\end{tikzpicture}
\caption{Use case diagram. The End User is the primary actor; the Processing Server (\texttt{server/main.py}) and the Cloud GPU (RunPod pod, or the local GPU) are secondary actors. Take Measurement currently runs in the desktop viewer; the web version is FYP-2 work.}
\label{fig:usecase}
\end{figure}

\section{Detailed Use Cases}

\paragraph{UC-1: Record Indoor Scan}
\begin{itemize}
    \item \textbf{Actor}: End User
    \item \textbf{Precondition}: App installed on ARCore-compatible Android device
    \item \textbf{Main flow}: User opens app $\rightarrow$ taps Record $\rightarrow$ walks through room $\rightarrow$ receives live IQA feedback (blur, speed, darkness, texture warnings) and coverage overlay (red/yellow/green dots) $\rightarrow$ taps Stop $\rightarrow$ sees quality summary
    \item \textbf{Postcondition}: Dataset saved locally (video, poses, depth, quality, meta)
\end{itemize}

\paragraph{UC-2: Upload and Process}
\begin{itemize}
    \item \textbf{Actor}: End User
    \item \textbf{Precondition}: Recording completed; server URL configured
    \item \textbf{Main flow}: User taps ``Upload \& Process'' $\rightarrow$ app zips depth folder, uploads all files $\rightarrow$ server creates job, dispatches to GPU $\rightarrow$ app polls progress every 3\,s showing step, percentage, logs $\rightarrow$ completion or failure reported
    \item \textbf{Postcondition}: 3D model (splat.ply) available for download
    \item \textbf{Exception}: Upload fails $\rightarrow$ retry button; pipeline step fails $\rightarrow$ error message with step name
\end{itemize}

\section{Functional Requirements}
Table~\ref{tab:fr} lists the functional requirements by module. Each has a priority, a status, a verifiable acceptance check, and its trace to the implementing file and the test case (Chapter~\ref{ch:testing}) that checks it.

{\footnotesize\setlength{\tabcolsep}{3pt}
\begin{longtable}{p{0.07\textwidth}p{0.27\textwidth}p{0.06\textwidth}p{0.10\textwidth}p{0.24\textwidth}p{0.17\textwidth}}
\caption{Functional requirements with acceptance checks and traceability}\label{tab:fr}\\
\toprule
\textbf{ID} & \textbf{Requirement (dependency)} & \textbf{Pri.} & \textbf{Status} & \textbf{Acceptance check} & \textbf{Implementation; test} \\
\midrule
\endfirsthead
\toprule
\textbf{ID} & \textbf{Requirement (dependency)} & \textbf{Pri.} & \textbf{Status} & \textbf{Acceptance check} & \textbf{Implementation; test} \\
\midrule
\endhead
\multicolumn{6}{l}{\textit{Mobile capture module}} \\
FR-M1 & Record video at 30\,FPS with per-frame camera poses (position, quaternion, intrinsics) in \texttt{poses.csv} (ARCore SDK) & Must & Done & One pose row per tracked video frame & \texttt{CaptureLogger.java}; TC-1 \\
FR-M2 & Log ARCore tracking state and failure reasons in \texttt{tracking.csv} (FR-M1) & Must & Done & One row per frame with state and reason & \texttt{CaptureLogger.java}; TC-1 \\
FR-M3 & Save ARCore feature points with confidence in \texttt{points.csv} (ARCore SDK) & Should & Done & Points written with timestamps & \texttt{CaptureLogger.java}; TC-1 \\
FR-M4 & Save raw depth maps with confidence at 5\,Hz in \texttt{depth/} (ARCore Depth API) & Should & Done & Depth written at most every 0.2\,s, only when ARCore produced a new map & \texttt{CaptureLogger.java}; TC-1 \\
FR-M5 & Compute and display live IQA: blur, brightness, motion speed, texture (FR-M1) & Must & Done & Warning shown on blurred, dark, fast or textureless frames & \texttt{QualityMonitor.java}; TC-2 \\
FR-M6 & Show a coverage overlay: red, yellow, green by observed directions (FR-M4) & Must & Done & Dot colour follows the 1 / 2--3 / $\geq$4 direction rule & \texttt{CoverageMap.java}; TC-3 \\
FR-M7 & Upload the complete dataset by multipart HTTP POST (server available) & Must & Done & Server returns a job ID; files saved & \texttt{ProcessingActivity.java}; TC-8 \\
FR-M8 & Configure the server URL from the settings menu & Must & Done & URL persists across app restarts & \texttt{HelloArActivity.java}; -- \\
FR-M9 & Show progress: step, percentage, elapsed time, recent log lines (FR-S5) & Must & Done & Fields update from \texttt{/status} every 3\,s & \texttt{ProcessingActivity.java}; TC-9 \\
\multicolumn{6}{l}{\textit{Processing server module}} \\
FR-S1 & Accept multipart uploads (video, poses, tracking, points, quality, meta, \texttt{depth.zip}) & Must & Done & Missing file rejected; corrupt zip fails cleanly & \texttt{main.py}; TC-8 \\
FR-S2 & Create a unique job ID per upload and track its status & Must & Done & Status has the fields the app reads & \texttt{pipeline.py} (\texttt{Job}); TC-9 \\
FR-S3 & Run the pipeline with per-step timeout watchdogs (pipeline scripts) & Must & Done & Step over its timeout is killed; job failed & \texttt{pipeline.py} (\texttt{StepConfig}); TC-15, TC-14 \\
FR-S4 & Support local GPU and RunPod cloud GPU modes (RunPod API key) & Must & Done & Mode switch via \texttt{/config} reaches the right starter & \texttt{runpod\_client.py}; TC-7, TC-10 \\
FR-S5 & Expose a progress-polling endpoint (FR-S2) & Must & Done & \texttt{GET /status/\{id\}}; unknown job 404 & \texttt{main.py}; TC-9 \\
FR-S6 & Provide a download endpoint for \texttt{splat.ply} (FR-S3) & Must & Done & 404 unknown, 409 unfinished, file when done & \texttt{main.py}; TC-16 \\
FR-S7 & In RunPod mode, always terminate the pod after processing, even on failure (FR-S4) & Must & Done & \texttt{terminate\_pod()} on every path where a pod exists & \texttt{runpod\_client.py}; TC-7 \\
\multicolumn{6}{l}{\textit{Reconstruction pipeline module}} \\
FR-P1 & Extract keyframes by gyro-aware sharpness prediction (exposure metadata) & Must & Partial: \texttt{select\_sharp.py} validated (v4--v6), not yet in the one-command pipeline & Least predicted blur per 0.2\,s window & \texttt{prepare.py}, \texttt{select\_sharp.py}; TC-5 \\
FR-P2 & Run COLMAP SfM with SIFT, lens self-calibration and global mapping (COLMAP 4.2) & Must & Done & Frames register; reprojection error $<$1\,px & \texttt{sfm\_colmap.py}; TC-5, TC-18 \\
FR-P3 & Compute a Sim3 alignment from COLMAP to ARCore's metric frame (FR-P2, FR-M1) & Must & Done & Model distances within 5\% of tape measure; measured mean 0.19\%, worst 0.29\% (Table~\ref{tab:tape}) & \texttt{sfm\_colmap.py}; TC-6 \\
FR-P4 & Train splatfacto and export \texttt{splat.ply} (Nerfstudio 1.1.5) & Must & Done & Held-out PSNR $\geq$ 25\,dB, SSIM $\geq$ 0.85; v6: 26.52\,dB / 0.8886 & \texttt{pipeline.py}; TC-13 \\
FR-P5 & Bridge ARCore tracking gaps with the IMU (FR-M2) & Should & Done & Bridged poses land on both gap anchors & \texttt{fallback.py}; TC-4 \\
\bottomrule
\end{longtable}
}

\section{Non-Functional Requirements}
Table~\ref{tab:nfr} lists the non-functional requirements with their measurement method and current status. Several are not yet measured; they are marked as such rather than assumed.

{\footnotesize\setlength{\tabcolsep}{3pt}
\begin{longtable}{p{0.07\textwidth}p{0.10\textwidth}p{0.33\textwidth}p{0.24\textwidth}p{0.18\textwidth}}
\caption{Non-functional requirements and how each is measured}\label{tab:nfr}\\
\toprule
\textbf{ID} & \textbf{Category} & \textbf{Requirement} & \textbf{Measurement} & \textbf{Status} \\
\midrule
\endfirsthead
\toprule
\textbf{ID} & \textbf{Category} & \textbf{Requirement} & \textbf{Measurement} & \textbf{Status} \\
\midrule
\endhead
NFR-R1 & Reliability & The cloud GPU pod shall be terminated within 60\,s of job completion or failure (prevents billing leaks). & Server log timestamps on a live pod & Not yet measured \\
NFR-R2 & Reliability & Each pipeline step shall have a timeout watchdog; if exceeded, the step is killed and the job marked failed. & Timeouts in \texttt{StepConfig} (TC-15); kill path (TC-14) & Implemented; kill path not yet executed \\
NFR-U1 & Usability & The capture app shall give real-time visual feedback without the user needing to understand 3D reconstruction. & User testing with external users & Not yet executed (planned FYP-2) \\
NFR-U2 & Usability & Processing progress shall show a progress bar, step name and elapsed time. & UI inspection & Implemented \\
NFR-P1 & Performance & End-to-end processing (upload to exported \texttt{splat.ply}) shall take under 35\,min on an NVIDIA RTX 4090 cloud GPU. FYP-2 target and a projection; the only measured time is $\sim$75\,min on the RTX 3070 (matching 23, mapping 6, training 45\,min). & Wall-clock time on an RTX 4090 pod & \textit{Outstanding: not yet measured} \\
NFR-P2 & Performance & The app shall keep a 30\,FPS camera preview while recording (no frame drops from logging). & No frame drops in logcat over a 3-minute recording & Not yet measured \\
NFR-P3 & Performance & The recording fast-path shall skip non-essential rendering (planes, virtual objects, PBR lighting). & Code inspection; GL thread profiling & Implemented; not profiled \\
NFR-S1 & Security & The RunPod API key shall be loaded from environment variables or a \texttt{.env} file, never hardcoded. & Code review; CI secrets check; \texttt{/config} never returns the key (TC-10) & Implemented \\
NFR-S2 & Security & The \texttt{.env} file shall be excluded from version control via \texttt{.gitignore}. & \texttt{.gitignore} inspection & \textbf{Not met}: \texttt{.env} is tracked in the repository \\
\bottomrule
\end{longtable}
}

\section{Acceptance Criteria}
\label{sec:acceptance}
The acceptance check of every requirement is in Tables~\ref{tab:fr} and~\ref{tab:nfr}. The project-level acceptance thresholds are: held-out PSNR $\geq$ 25\,dB and SSIM $\geq$ 0.85 (FR-P4; met by v6), distances within 5\% of tape measure (FR-P3; met, worst case 0.29\%), and processing under 35\,min on an RTX 4090 (NFR-P1; not yet measured).
